% Generated by report/tools/tables.py from the result files of the repository. Do not edit.
\begin{table}[tbp]
\centering
\footnotesize
\setlength{\tabcolsep}{3.5pt}
\begin{threeparttable}
\caption[Trajectoire 2018 à 2026]{Trajectoire du RL classique de 2018 à 2026 : part des articles qui remplissent chaque critère, en pourcentage, par année. \textbf{En gras}, ce qui passe la correction de Holm. Source : \file{study/trajectory.json} (\file{study/trajectory.py}).}
\label{tab:trajectoire}
\begin{tabularx}{\linewidth}{@{}l>{\raggedright\arraybackslash}Xrrrrrrrr@{}}
\toprule
Code & Critère & 2018 & 2019 & 2020 & 2022 & 2024 & 2026 & $p$ tendance & $p$ Holm \\
 & $n$ & 48 & 20 & 20 & 20 & 20 & 50 & &  \\
\midrule
\textbf{Q5} & \textbf{Code en lien} & \textbf{31} & \textbf{60} & \textbf{50} & \textbf{70} & \textbf{75} & \textbf{74} & \textbf{2\,×\,10\textsuperscript{−5}} & \textbf{7\,×\,10\textsuperscript{−5}} \\
\textbf{Y9} & \textbf{Sensibilité à un hyperparamètre} & \textbf{44} & \textbf{35} & \textbf{45} & \textbf{60} & \textbf{85} & \textbf{80} & \textbf{2\,×\,10\textsuperscript{−6}} & \textbf{9\,×\,10\textsuperscript{−6}} \\
\textbf{X7} & \textbf{Calcul chiffré} & \textbf{27} & \textbf{30} & \textbf{40} & \textbf{40} & \textbf{55} & \textbf{58} & \textbf{0,0006} & \textbf{0,002} \\
\textbf{X2} & \textbf{Incertitude rapportée} & \textbf{63} & \textbf{80} & \textbf{85} & \textbf{95} & \textbf{100} & \textbf{88} & \textbf{0,0007} & \textbf{0,002} \\
Q8 & Test statistique appliqué & 6 & 5 & 10 & 10 & 25 & 12 & 0,13 & 0,13 \\
Q3 & Tous les hyperparamètres de la méthode & 55\,{\scriptsize (47)} & 45 & 55 & 70 & 55 & 72 &  &  \\
Q4 & Tous les hyperparamètres des références & 35 & 16\,{\scriptsize (19)} & 45 & 25 & 40 & 29\,{\scriptsize (49)} &  &  \\
Q6 & Choix des hyperparamètres décrit & 26\,{\scriptsize (47)} & 30 & 40 & 40 & 35 & 48 &  &  \\
Q7 & Évaluation sur conditions tenues à part & 33 & 25 & 20 & 20 & 55 & 45\,{\scriptsize (49)} &  &  \\
X4 & Références réglées avec le même soin & 13 & 26\,{\scriptsize (19)} & 35 & 15 & 25 & 16\,{\scriptsize (49)} &  &  \\
Y7 & Version des bancs d'essai & 23 & 25 & 20 & 40 & 20 & 26 &  &  \\
\midrule
\multicolumn{10}{@{}l}{\emph{Différences entre périodes (Newcombe, Fisher, Holm sur les cinq critères testés)}} \\
Code & Critère & \multicolumn{3}{r}{\makecell[r]{2019-2020 contre 2018\\{}différence [IC]}} & $p$ Holm & \multicolumn{3}{r}{\makecell[r]{2022-2024 contre 2019-2020\\{}différence [IC]}} & $p$ Holm \\
\cmidrule(lr){1-10}
Q5 & Code en lien & \multicolumn{3}{r}{+24 [+3\,;\,+42]} & 0,16 & \multicolumn{3}{r}{+17 [−3\,;\,+36]} & 0,49 \\
Y9 & Sensibilité à un hyperparamètre & \multicolumn{3}{r}{−4 [−23\,;\,+16]} & 1 & \multicolumn{3}{r}{\textbf{+32 [+11\,;\,+50]}} & \textbf{0,03} \\
X7 & Calcul chiffré & \multicolumn{3}{r}{+8 [−11\,;\,+27]} & 1 & \multicolumn{3}{r}{+13 [−9\,;\,+32]} & 0,62 \\
X2 & Incertitude rapportée & \multicolumn{3}{r}{+20 [+1\,;\,+37]} & 0,23 & \multicolumn{3}{r}{+15 [+1\,;\,+30]} & 0,23 \\
Q8 & Test statistique appliqué & \multicolumn{3}{r}{+1 [−10\,;\,+14]} & 1 & \multicolumn{3}{r}{+10 [−5\,;\,+25]} & 0,62 \\
\bottomrule
\end{tabularx}
\begin{tablenotes}[flushleft]\scriptsize
\item[] 2018 : notre lecture des articles à expériences ; 2019, 2020, 2022 et 2024 : les 20 premiers articles expérimentaux de RL classique par année, tirés au hasard dans ICML et NeurIPS (ICLR absente, son site refusant les téléchargements automatiques) ; 2026 : l'échantillon classique de la comparaison principale. Le nombre entre parenthèses donne le dénominateur quand il diffère de l'effectif de la colonne (critère sans objet pour certains articles). Avec 20 articles, une part a une incertitude d'environ vingt points de part et d'autre. Pour les cinq critères testés (Q5, Y9, X7, X2, Q8) : test de tendance de Cochran-Armitage sur les six années (scores égaux aux années), $p$ bilatéral et $p$ corrigé par Holm sur les cinq ; les six autres lignes sont descriptives. Partie basse : différence immédiate (2019 et 2020 réunis contre 2018) et différence tardive (2022 et 2024 contre 2019 et 2020), en points, intervalle de Newcombe à 95~\%, Holm sur les cinq critères dans chaque famille.
\end{tablenotes}
\end{threeparttable}
\end{table}
